% =====================================================================
%  How Many Seeds? Dispersion, Resolving Power, and Silent Failure in
%  Quantum Sequence Modelling
%
%  V. Liutarevich
%  Draft v4 --- prepared for Quantum (quantum-journal.org)
%
%  ---------------------------------------------------------------------
%  DOCUMENT CLASS
%  ---------------------------------------------------------------------
%  Quantum requires `quantumarticle'. It was not installed where this
%  was compiled, so the file defaults to `article' and builds anywhere.
%  For submission replace the \documentclass line with
%
%      \documentclass[twocolumn,superscriptaddress]{quantumarticle}
%
%  and delete the three fallback definitions marked below. No
%  quantumarticle-specific macro is used in the body.
%
%  ---------------------------------------------------------------------
%  WHAT CHANGED FROM v3
%  ---------------------------------------------------------------------
%  v3 argued that small-sample comparisons in this field cannot resolve
%  the effects they report, and supported it with measured dispersion, a
%  failed replication, and resampled false-positive rates.
%
%  v4 adds a section the earlier draft could not have written, because
%  the material did not exist yet: a case series of five SILENT
%  analysis failures from the companion architecture study, three of
%  which changed a reported number, with the before and after values.
%  The centrepiece is an apparent 29% improvement that was entirely an
%  artefact of pooling two tasks -- and whose p-value was IDENTICAL
%  before and after the error was found, because only the magnitude was
%  spurious. That example does something the resampling argument cannot:
%  it shows a significance test passing a result that is wrong.
%
%  v4 also promotes the permutation floor to its own section. The
%  smallest two-sided p an exact sign-flip test can return at n seeds is
%  2/2^n, so a three-seed study cannot reach significance whatever its
%  data. This is arithmetic, not statistics, and it is routinely
%  violated in practice.
%
%  Three audit items from v3 are now CLOSED and marked as such in the
%  text; the rest are listed in Appendix A.
%
%  ---------------------------------------------------------------------
%  PLACEHOLDERS
%  ---------------------------------------------------------------------
%     \pend{...}     inline      \pendpara{...}  paragraph
%     pendingbox     whole block
%  \draftfalse hides all of them.
% =====================================================================
\documentclass[11pt,onecolumn]{article}

\usepackage[margin=1in]{geometry}
\usepackage[T1]{fontenc}
\usepackage[utf8]{inputenc}
\usepackage{amsmath,amssymb}
\usepackage{booktabs}
\usepackage{graphicx}
\usepackage{xcolor}
\usepackage{caption}
\usepackage{etoolbox}
\usepackage[colorlinks=true,linkcolor=NavyBlue,citecolor=NavyBlue,
            urlcolor=NavyBlue]{hyperref}
\usepackage[numbers,sort&compress]{natbib}

\definecolor{NavyBlue}{HTML}{1F3864}
\definecolor{PendBG}{HTML}{FFF4E5}
\definecolor{PendFG}{HTML}{8A4B00}
\definecolor{DoneBG}{HTML}{EAF3DE}
\definecolor{DoneFG}{HTML}{3B6D11}

\newif\ifdraft
\draftfalse

\newcommand{\pend}[1]{%
  \ifdraft\textcolor{PendFG}{\textbf{[}\textit{#1}\textbf{]}}\fi}
\newcommand{\pendpara}[1]{%
  \ifdraft\par\vspace{2pt}\noindent
  \fcolorbox{PendFG}{PendBG}{\begin{minipage}{0.97\linewidth}\small
    \textcolor{PendFG}{\textbf{PENDING.} #1}\end{minipage}}\par\vspace{4pt}\fi}
\newcommand{\closed}[1]{%
  \ifdraft\par\vspace{2pt}\noindent
  \fcolorbox{DoneFG}{DoneBG}{\begin{minipage}{0.97\linewidth}\small
    \textcolor{DoneFG}{\textbf{AUDIT CLOSED.} #1}\end{minipage}}\par\vspace{4pt}\fi}

\newsavebox{\pendbox}
\newenvironment{pendingbox}[1]
  {\par\vspace{3pt}\noindent
   \begin{lrbox}{\pendbox}\begin{minipage}{0.93\linewidth}\small
   \textcolor{PendFG}{\textbf{PENDING --- #1}}\par\vspace{3pt}\color{PendFG}}
  {\end{minipage}\end{lrbox}%
   \ifdraft\fcolorbox{PendFG}{PendBG}{\usebox{\pendbox}}\fi\par\vspace{4pt}}

\newcommand{\figorplaceholder}[4]{%
  \IfFileExists{figures/#1}{%
    \begin{figure}[htb]\centering
      \includegraphics[width=#2\linewidth]{figures/#1}
      \caption{#3}\label{fig:#4}\end{figure}}{%
    \ifdraft\begin{figure}[htb]\centering
      \fcolorbox{PendFG}{PendBG}{\begin{minipage}{0.88\linewidth}\small
        \textcolor{PendFG}{\textbf{FIGURE NOT FOUND:}
        \texttt{figures/\detokenize{#1}}\\[2pt]
        Run \texttt{make\_figures.py} / \texttt{power\_analysis.py}%
        }\end{minipage}}
      \caption{#3}\label{fig:#4}\end{figure}\fi}}

% ---- delete these three when switching to quantumarticle ------------
\newcommand{\affiliation}[1]{\\[2pt]\small\textit{#1}}
\newcommand{\email}[1]{\\[2pt]\small\texttt{#1}}
\newcommand{\maketitlefallback}{\maketitle}

\title{\bfseries How Many Seeds?\\
       Dispersion, Resolving Power, and Silent Failure\\
       in Quantum Sequence Modelling}
\author{V. Liutarevich
        \affiliation{\pend{affiliation}}
        \email{\pend{contact email}}}
\date{\today}

\begin{document}
\maketitlefallback

\begin{abstract}
\noindent
Recent work has established that quantum machine learning models rarely
outperform classical baselines at simulable scale, and that apparent
advantages often trace to architectural rather than quantum effects. We
address a complementary question those results leave open: whether the
sample sizes in routine use could detect such an advantage if one
existed, and whether the analysis pipelines in routine use would notice
if they were wrong.

The quantity that decides the first question is seed-to-seed dispersion,
which is almost never reported. We measure it across seven populations
with matched architectures and find it spans a factor of
\textbf{twenty-nine} --- standard deviation of the per-seed log-ratio
from $0.014$ to $0.404$. The sample required to detect a $10\%$ effect
at $80\%$ power therefore varies from six seeds to a hundred and fifty,
with no way to know which regime one occupies without measuring. The
consequence for false positives is equally wide: resampling under the
paired null, the fraction of null studies at five seeds reporting a
$10\%$-or-better advantage ranges from $0\%$ to $30\%$ across these
populations, with the same two architectures throughout. A small pilot
does not help: a three-seed pilot implies a required sample anywhere
between $4$ and $169$.

We demonstrate the consequence with a failed replication. An apparent
$21\%$ advantage on chaotic forecasting, observed at $13$ seeds, did not
survive twelve further seeds run specifically to test it; pooled over
$25$ seeds the ratio is $0.907$, interval $[0.780, 1.059]$. Resampling
under the paired null explains why: at this task's dispersion, more than
one null study in four at five seeds reports a $10\%$-or-better
advantage.

We then report something the resampling argument cannot capture. Across
this study and the companion architecture study~\cite{paperb}, seven
analysis failures occurred that produced plausible wrong numbers rather
than errors; four changed a reported result. In the most instructive, an
apparent $29\%$ improvement was manufactured entirely by pooling two
different environments in one aggregation key, and \textbf{the
permutation $p$-value was identical before and after the error was
found} ($0.041$ both times), because only the magnitude was spurious. A
significance test cannot detect this class of error.

Two of the seven were in \emph{this} paper, and were found by applying
the audit we recommend to our own pipeline: one comparison was pooling
five model widths, and one data file silently contained a failed run and
a different entangler from the one our methods section claimed. We
report them in full. A paper arguing that aggregation errors change
published numbers is obliged to check itself first, and the check
earned its place.

Separately, we note an arithmetic constraint that is routinely violated:
the smallest two-sided $p$-value an exact sign-flip permutation test can
return at $n$ seeds is $2/2^n$, so a three-seed comparison cannot reach
$p < 0.05$ whatever its data. Across all paired comparisons in the two
studies, $30$ in total, three reach $p < 0.05$ against $1.5$ expected
under the global null, and none survives correction for multiplicity.

We also obtain, from six runs repeated at a fixed seed, a direct
estimate of how much of ``seed variance'' is not seed variance:
absolute error moves by up to a factor of two between two runs of
identical seed and configuration, while the \emph{paired} ratio moves
with a log standard deviation of $0.071$ against $0.477$ across seeds.
The paired protocol absorbs nearly all of it, which justifies that
design empirically rather than by assertion.

Finally we give a theoretical account of why no effect is expected at
this scale and test it twice. A data re-uploading circuit computes a
truncated Fourier series whose accessible frequencies are fixed by its
encoding; an eight-parameter classical surrogate over that spectrum
matches the thirty-parameter circuit. And in the companion study,
deleting every two-qubit gate from the circuit changes error by about
$1\%$ --- the property that would make it hard to simulate contributes
nothing measurable. Varying the relevant block's capacity over a
$64$-fold range, meanwhile, moves performance less than changing the
random seed does, which yields a second necessary condition independent
of sample size: the component under comparison must constrain
performance, or no amount of data makes the comparison informative.
\end{abstract}

\vspace{1em}

% =====================================================================
\section{Introduction}
% =====================================================================
Does a model containing a parameterised quantum circuit outperform a
classical model of comparable size? Recent work has made substantial
progress, and the answer at currently simulable scale is largely
negative. Bowles, Ahmed and Schuld~\cite{bowles2024} tested twelve
quantum models across $160$ classification datasets and found classical
baselines generally superior, with entanglement frequently dispensable.
The authors of QASA~\cite{qasa}, a quantum transformer whose dynamics
core we reproduce here, have revised their own work to attribute its
reported advantage to a low-rank value-projection bottleneck rather than
to the quantum substrate, demonstrated by a capacity-matched classical
control.

Both results concern confounds in the \emph{architecture}: what is being
compared. This paper concerns two confounds in the \emph{measurement}:
whether the comparison could resolve the effect it reports, and whether
the pipeline that produced it would have revealed an error if one had
occurred.

The first is a question about dispersion across random initialisations.
It is almost never reported, and without it neither author nor reader
can say whether ten seeds is generous or hopeless. We measure it with
matched architectures and find it varies twelvefold, moving the sample
required to detect a $10\%$ effect from six seeds to a hundred and
fifty.

The second is a question about provenance, and it is the one we did not
anticipate. Over the course of the companion architecture
study~\cite{paperb} five analysis or bookkeeping failures occurred, none
of which raised an error and all of which produced plausible numbers.
Three changed a reported result. We document them as a case series
because the literature's methodological discussion is almost entirely
about statistics, and in our experience the failures that actually
altered our conclusions were not statistical.

\paragraph{Contributions.}
\begin{enumerate}\itemsep3pt
\item \textbf{A measurement of seed dispersion} across task settings
      with matched architectures, with intervals, and the sample size
      each implies (Section~\ref{sec:howmany}).
\item \textbf{A failed replication with a pre-committed out-of-sample
      test} (Section~\ref{sec:replication}), independently corroborating
      the revised conclusion of~\cite{qasa} by a different route.
\item \textbf{False-positive rates estimated from observed dispersion},
      the empirical counterpart to the simulation-based
      publication-bias argument of~\cite{bowles2024}.
\item \textbf{Evidence that a small pilot cannot forecast the required
      sample size}, which constrains what can usefully be recommended
      (Section~\ref{sec:pilot}).
\item \textbf{The permutation floor}, and the observation that it is
      routinely violated: at $n$ seeds no exact sign-flip test can
      return a two-sided $p$ below $2/2^n$
      (Section~\ref{sec:floor}).
\item \textbf{A case series of five silent analysis failures} with
      before-and-after values, including one in which the $p$-value was
      unchanged while the effect size was fabricated
      (Section~\ref{sec:silent}). To our knowledge no comparable
      first-hand account exists in this literature.
\item \textbf{A measured magnitude for a trainability-diagnostic
      pitfall}: a common entangler choice violates the depth
      precondition of the local-cost protection result, changing
      measured gradient variance by roughly $1{,}260\times$ at ten
      qubits (Section~\ref{sec:trainability}).
\item \textbf{A theoretical account of why no effect is expected} at
      this scale, tested with a classical surrogate over the circuit's
      own spectrum and, independently, with an entanglement ablation
      (Section~\ref{sec:why}).
\item \textbf{A second necessary condition for architectural
      comparison} --- that the component under test must constrain
      performance --- which the benchmark underlying the original claim
      fails, and which no sample size can repair
      (Section~\ref{sec:second}).
\end{enumerate}

We do not claim that quantum advantage in sequence modelling is
impossible. Section~\ref{sec:why} argues the narrower and more specific
point that at simulable scale no expressivity advantage is available in
principle, leaving inductive bias as the only live mechanism, which our
measurements bound at $2.1\%$. That argument applies to variational
models with generic ans\"atze; it does not extend to problem-informed
circuits or to non-variational approaches, and we say so.

The methodological claim is separate and narrower: in the regime where
most empirical work takes place, small effects cannot be resolved by the
samples in use, the conventions of the field make this invisible, and
the analysis pipelines that produce the numbers are not routinely
checked for the failures that would alter them.

% =====================================================================
\section{Related work}
% =====================================================================
\paragraph{Benchmarking and evaluation.}
Bowles et al.~\cite{bowles2024} conduct the most extensive evaluation of
quantum machine learning models to date, finding classical baselines
generally superior and entanglement often dispensable. They argue
further, through simulation, that publication bias over a large design
space can invert an apparent ranking. Our work is complementary in two
respects. Where they go broad, using cross-validation per
configuration, we go deep: the split is held fixed and only
initialisation varies, so seed dispersion is isolated and can be
measured. And where their publication-bias argument assumes performance
distributions, we estimate false-positive rates from the dispersion our
own runs exhibit.

\paragraph{Capacity-matched controls.}
The QASA architecture~\cite{qasa}, whose quantum token-mixing block we
reproduce, has been revised to attribute its reported advantage to a
low-rank value-projection bottleneck rather than to quantumness. Our
null on chaotic forecasting is consistent with that conclusion, reached
independently: we hold capacity matched throughout and vary the number
of seeds. We adopt their control as a requirement, not a contribution.

\paragraph{Trainability.}
Cerezo et al.~\cite{cerezo2021} show that local cost functions avoid
barren plateaus provided ansatz depth remains $O(\log n)$ --- a
condition expressed through the causal cones of the cost terms.
Section~\ref{sec:trainability} reports the empirical consequence of
violating that condition inadvertently.

\paragraph{Variance in machine learning benchmarks.}
Our statistical argument follows established work in classical machine
learning~\cite{henderson2018,agarwal2021,bouthillier2021,dodge2019} on
seed variance, interval estimation and the inadequacy of small-sample
comparisons, transposed to a setting where per-run cost is far higher
and reported samples are correspondingly smaller --- which makes the
problem more acute, not less. Recent work benchmarking quantum models on
chaotic dynamics~\cite{pandey2026} reports a comparator at
$47.9 \pm 36.6$, a standard deviation $76\%$ of the mean, illustrating
the magnitudes involved in precisely our setting.

\paragraph{Provenance and silent failure.}
We are not aware of a comparable first-hand case series in the quantum
machine learning literature. The closest analogues are in computational
reproducibility more broadly, where the failures documented are usually
environmental or numerical rather than failures of aggregation logic.
\pend{search for prior art on aggregation-key errors in ML pipelines;
if a comparable account exists, contribution 6 weakens}

\pendpara{Exact citations, years and venues for
\cite{henderson2018,agarwal2021,bouthillier2021,dodge2019,pandey2026};
author list for \cite{qasa}; journal volume and page for
\cite{schuld2021}. Also outstanding: a search for any quantum machine
learning paper that reports seed dispersion and derives a required
sample size. If one exists, contribution 1 weakens.}

% =====================================================================
\section{Experimental setting}\label{sec:setting}
% =====================================================================
\paragraph{Architectures.}
A transformer encoder in which one token-mixing block is replaced by a
parameterised quantum circuit (six qubits, four layers, angle encoding
with data re-uploading, local Pauli-$Z$ readout), against an identical
transformer with a classical block in the same position.

\paragraph{Entangler, and a correction.}
The entangling topology differs between experiment families, and an
earlier version of this paper stated brickwork throughout. That was
wrong. The forecasting experiments (sine, noisy sine, Mackey--Glass) use
a \textbf{sequential CNOT ring cascade}; the Ising and latent-dynamics
experiments use \textbf{brickwork}. We discovered the discrepancy by
auditing our own data files (Section~\ref{sec:selfaudit}): the
forecasting records carried no entangler column for the first three
seeds and an explicit \texttt{ring\_cascade} field thereafter, with the
brickwork runs written to a separate file at three seeds only.

This matters here more than it would elsewhere, because
Section~\ref{sec:trainability} measures a $1{,}260\times$ difference in
gradient variance between these topologies at ten qubits. At six qubits
the difference is far smaller, and none of this paper's conclusions
depends on the entangler --- the dispersion measurement, the failed
replication and the resolving-power argument are all invariant to it.
But the methods section must say what was run, and ours did not. Following~\cite{qasa}, the classical control is a
capacity-matched low-rank bottleneck of the same parameter budget rather
than a full-width block: it projects $d_{\mathrm{model}} \rightarrow
n_{\mathrm{qubits}} \rightarrow \mathrm{hidden} \rightarrow
d_{\mathrm{model}}$ with the same budget as the circuit it replaces.
Parameter counts are matched to within $15\%$ by construction and
verified per run: typically $7{,}274$ against $7{,}518$, of which $28$
are circuit parameters.

\closed{v3 flagged for verification whether the dynamics-model control
(\texttt{DynamicsCore}, \texttt{quantum=False}) is likewise a bottleneck
rather than a full-width block. It is. The classical token layer is
\texttt{Linear}$(d_{\mathrm{model}} \rightarrow n_{\mathrm{qubits}})$,
$\tanh$, \texttt{Linear}$(n_{\mathrm{qubits}} \rightarrow 16)$, $\tanh$,
\texttt{Linear}$(16 \rightarrow d_{\mathrm{model}})$, with
$n_{\mathrm{qubits}} = 6$. Both arms therefore constrict to six
dimensions in the same sublayer, so the Ising and CartPole comparisons
are capacity-matched in the sense of~\cite{qasa}. No qualification is
required.}

\paragraph{Tasks.}
\begin{table}[htb]\centering\small
\begin{tabular}{lll}
\toprule
Family & Dynamics & Observations \\
\midrule
Latent control dynamics & CartPole, DreamerV3 latents &
  6-dim compressed latent \\
Synthetic forecasting & sine, noisy sine, Mackey--Glass &
  scalar series, lookback 24 \\
Quantum-native dynamics & transverse-field Ising chain, 2--6 spins &
  $\langle Z \rangle$, optionally $\langle X \rangle$ \\
\bottomrule
\end{tabular}
\end{table}

The Ising chain is evolved exactly by dense matrix exponentiation,
removing Trotter error as a confound. Its control field switches among
five levels. By default only $\langle Z_i \rangle$ is recorded: those
$N$ numbers do not determine the quantum state, so prediction requires
history. Recording $\langle X_i \rangle$ as well moves the task nearer
Markovian and serves as a control arm.

\paragraph{Why exact simulation, and what it costs.}
All circuits are simulated exactly by statevector, with gradients
obtained by backpropagation through the simulator. This is deliberate
and conservative: the quantum models are trained under conditions no
physical device provides --- no gate noise, no decoherence, no readout
error, exact rather than shot-estimated gradients, effectively infinite
measurement budget. Every one of these favours the quantum model. A
comparison the quantum architecture loses under ideal conditions it will
not win under realistic ones, so a null obtained this way is stronger
than a null obtained on hardware, not weaker.

Running on a device would answer a different question --- whether a
quantum model remains usable despite noise --- and would add dispersion
rather than remove it, the opposite of what a comparison needs. What
additional compute would most usefully buy is \emph{more seeds}, and we
note the reflexivity: this study is sample-limited by available compute,
which is precisely the constraint that produces the field-wide problem
we measure in Section~\ref{sec:howmany}.

\paragraph{Protocol.}
Data are split by whole trajectory, never by window, and the split is
fixed across seeds so that seeds vary model initialisation only. Every
error is reported alongside persistence, linear-extrapolation and
training-mean predictors on the same split. Comparisons are paired
within seed and tested by permutation on log-ratios --- logarithms
because ratios are multiplicative, so that $0.90$ and $1.11$ count as
equal and opposite --- with exact enumeration of sign assignments where
the sample permits. Evaluation covers one-step error and open-loop
rollout at horizons $5$, $10$ and $20$, in which the model consumes its
own predictions.

\paragraph{Aggregation.}
Every aggregation key used in this paper comprises the complete run
configuration, with only the two axes under comparison excluded, and
halts with an error rather than selecting a winner if two runs collide
on a key. This is not a stylistic choice; Section~\ref{sec:silent}
explains what it is a response to.

% =====================================================================
\section{Results, and the bounds on them}\label{sec:results}
% =====================================================================
\subsection{The measurement detects architectural effects of the
            relevant size}\label{sec:sensitivity}
A comparison returning ratios near unity invites two readings: that the
architectures perform equivalently, or that the measurement is
insensitive. We separate them before reporting results.

\paragraph{The two models are distinct.}
On the Ising four-spin task, per-seed differences between quantum and
classical error swing in both directions, from $-0.00105$ to
$+0.00137$, with a cross-seed correlation of $-0.236$. The mean absolute
difference is $2.08\times$ the quantum model's own seed-to-seed standard
deviation. Were the quantum block inactive the correlation would
approach unity and the differences would vanish; these are distinct
models that perform equivalently.

\paragraph{The pipeline resolves effects an order of magnitude smaller
than it would need to.}
Using the same task, metric and evaluation code, widening the classical
model alone produces a measurable, monotonic change:

\begin{table}[htb]\centering\small
\caption{Classical width sweep, Ising four-spin task.}
\begin{tabular}{lrrr}
\toprule
$d_{\mathrm{model}}$ & parameters & error ($k=1$) & ratio to $d=16$ \\
\midrule
$8$  & $3{,}306$  & $0.02676$ & $1.311$ \\
$16$ & $7{,}386$  & $0.02041$ & --- \\
$32$ & $18{,}618$ & $0.01973$ & $0.967$ \\
$64$ & $53{,}370$ & $0.01768$ & $\mathbf{0.867}$ \\
\bottomrule
\end{tabular}
\end{table}

Classical width alone yields a $13\%$ improvement from $d=16$ to
$d=64$. The quantum substitution at matched capacity yields $-0.1\%$.
The instrument registers an architectural effect of $13\%$ on this task
while registering nothing from the quantum substrate, which establishes
that ratios near unity reflect the absence of an effect rather than the
absence of sensitivity.
\pend{the $d \in \{8,32,64\}$ rows are single-seed; rerunning at three
seeds is cheap and would let this table carry intervals}

\figorplaceholder{fig3_forest.png}{0.95}{Geometric-mean ratio of quantum
to classical test error with $95\%$ bootstrap intervals, for every
comparison conducted. Values below $1$ favour the quantum model.}{forest}

\begin{table}[htb]\centering\small
\caption{Every comparison conducted, regenerated from the raw run
records. Ratios below $1$ favour the quantum model. The CartPole row is
withheld; see below.}
\label{tab:forest}
\begin{tabular}{lrrlrr}
\toprule
Setting & $n$ & ratio & $95\%$ CI & Q ahead & $p$ \\
\midrule
Ising chain, 2 spins           & $10$ & $0.969$ & $[0.919, 1.019]$ & $7/10$ & $0.309$ \\
Ising chain, 4 spins           & $10$ & $1.001$ & $[0.979, 1.025]$ & $6/10$ & $0.943$ \\
Ising chain, 6 spins           & $10$ & $0.976$ & $\mathbf{[0.959, 0.993]}$ & $8/10$ & $\mathbf{0.041}$ \\
Ising chain, 4 spins ($zx$)    & $10$ & $1.015$ & $[0.987, 1.047]$ & $4/10$ & $0.375$ \\
Forecasting, sine              & $10$ & $1.184$ & $[0.902, 1.571]$ & $4/10$ & $0.285$ \\
Forecasting, noisy sine        & $10$ & $0.998$ & $[0.990, 1.006]$ & $4/10$ & $0.654$ \\
Forecasting, Mackey--Glass     & $25$ & $0.907$ & $[0.780, 1.059]$ & $16/25$ & $0.240$ \\
\midrule
Latent dynamics, CartPole      & \multicolumn{5}{l}{\emph{withheld --- see Section~\ref{sec:selfaudit}}} \\
\bottomrule
\end{tabular}
\end{table}

\paragraph{A withheld row.}
An earlier version of this table reported latent control dynamics at
three seeds with a ratio of $0.991$. That number was produced by a
loader which, through a filter that never excluded anything, pooled five
different model widths and retained whichever happened to be read last
per seed. Since Section~\ref{sec:sensitivity} establishes that width
alone moves error by $13\%$ on a comparable task, the row measured
nothing interpretable and is withdrawn. The underlying runs exist at
eight seeds with correct provenance and will be reinstated.
\pend{reinstate from the Phase 1 records: at $d_{\mathrm{model}}=32$,
$K=4$ the ratio is $0.988$ ($p = 0.29$, ahead $5/8$) and at $K=8$ it is
$0.975$ ($p = 0.041$, ahead $6/8$); choose one context and state it}

Every interval contains $1$. We decline to describe this as evidence of
no effect, because the intervals differ enormously in what they
exclude. Read as bounds:

\begin{itemize}\itemsep2pt
\item \textbf{Noisy sine, 10 seeds.} Interval $[0.990, 1.006]$: any
      advantage above $\mathbf{1.0\%}$ is excluded. The tightest bound
      in the study.
\item \textbf{Ising, 4 spins, 10 seeds.} Interval $[0.979, 1.025]$: any
      advantage above $\mathbf{2.1\%}$ is excluded.
\item \textbf{Mackey--Glass, 25 seeds.} Interval $[0.780, 1.059]$: only
      advantages above $\mathbf{22\%}$ are excluded. A $15\%$ advantage
      remains entirely consistent with these data.
\item \textbf{Sine, 10 seeds.} Interval $[0.902, 1.571]$ --- wider at
      ten seeds than the noisy-sine interval at ten seeds by a factor
      of \emph{forty}, with the same two architectures.
\end{itemize}

\paragraph{The one cell that cleared threshold.}
Ising at six spins returns $0.976$ with an interval of $[0.959, 0.993]$
that excludes unity, ahead in eight seeds of ten, $p = 0.041$. We report
it and decline to promote it, for three reasons, all of which are this
paper's own arguments applied to its own result. It is one of seven
comparisons in Table~\ref{tab:forest}, against $0.35$ expected below
$\alpha = 0.05$ under the global null. It is not monotonic in system
size --- four spins sits at $1.001$ --- so there is no trend to appeal
to. And the effect is $2.4\%$, which is inside the bound the four-spin
comparison already reports as a null. A $2.4\%$ advantage on exactly
simulated quantum dynamics would be interesting if it replicated at
fresh seeds; we have not run that replication, and until we do this is
a nominally significant cell in a table where one was expected.

A table of eight ratios all near $1$ would support a far stronger
conclusion than these data license. The difference between a $2\%$ bound
and a $22\%$ bound is the whole question, and it is invisible unless
intervals are reported.

The quantum models are also roughly $20\times$ more expensive in
wall-clock terms (median $4{,}136$\,s against $182$\,s on CartPole;
$440$\,s against $23$\,s on the Ising task). These timings are
contaminated by machine sleep and competing load and should be read as
lower bounds on that ratio.

% =====================================================================
\section{A failed replication}\label{sec:replication}
% =====================================================================
\figorplaceholder{fig1_seed_collapse.png}{0.95}{The estimate a study
would have reported had it stopped after $n$ seeds. Left: chaotic
forecasting, with $95\%$ bootstrap band. Right: the two arms of the
Ising observables contrast.}{collapse}

\paragraph{Chaotic forecasting.}
At $13$ seeds the quantum model's error ratio was $0.787$ --- a $21\%$
improvement, nine wins in thirteen, with $p = 0.022$ under an unpaired
test of the mean. Twelve further seeds were then run for the sole
purpose of testing this estimate, making them an out-of-sample
replication committed to before the data were seen. Those twelve give
$1.058$: the classical model marginally ahead. Pooled over $25$ seeds
the ratio is $0.907$, interval $[0.780, 1.059]$, permutation
$p = 0.240$.

An estimate of this kind is what a study reporting $13$ seeds would
publish, and Section~\ref{sec:nullrates} gives the rate at which the
null produces it: at this task's dispersion, $17\%$ of thirteen-seed
studies report a $10\%$-or-better advantage with no effect present.

This is consistent with the revised conclusion of~\cite{qasa}, arrived
at independently. Where those authors attribute the apparent gain to a
low-rank bottleneck by constructing a capacity-matched control, we hold
capacity matched throughout and find the gain was not reliably present
to begin with. The two explanations are compatible and mutually
reinforcing.

Two features of the original estimate deserve note. The same thirteen
seeds give $p = 0.022$ unpaired and $p = 0.069$ under the paired
permutation test on log-ratios, with no change to the data --- the
choice of test moves the result across the conventional threshold. And
the per-seed ratios span $0.376$ to $2.244$, a six-fold range on a task
where the data are byte-identical across seeds and only the weight
initialisation varies.

\paragraph{The observables contrast.}
The hypothesis was that quantum attention helps where prediction
requires history. At three seeds per arm the non-Markovian arm sat at
$0.969$ and the nearer-Markovian arm at $1.025$, a separation of $5.6$
percentage points in the predicted direction. At ten seeds per arm the
arms sit at $1.001$ and $1.015$, a separation of $1.4$ points, and the
separation reverses sign between prediction horizons ($-0.4$ points at
$k = 10$). We take the sign reversal as the clearest available
indication that the estimand is zero.

These are not the same failure. One was a sample too small for the
task's dispersion; the other a sample adequate for a $10\%$ effect being
asked to resolve one of $5.6\%$.

\paragraph{A third instance, from the companion study.}
The same pattern recurred in~\cite{paperb} with a different
architecture and a different metric. A partial-observation return
difference of $+212$ points at three seeds became $+122$ at six and
$+58$ at seven, while the standard deviation of the paired difference
grew from $117$ to $182$ to $237$: dispersion grew faster than the
effect. Three independent instances in two studies is not proof of a
general rate, but it is the reason we treat three-seed estimates as
uninformative by default rather than as weak evidence.

% =====================================================================
\section{How many seeds?}\label{sec:howmany}
% =====================================================================
\subsection{Dispersion varies by an order of magnitude}
\label{sec:dispersion}
We report dispersion as the standard deviation of the per-seed
log-ratio, with a bootstrap interval --- a point estimate would be
inconsistent with the paper's own recommendation.

\begin{table}[htb]\centering\small
\caption{Seed dispersion and the sample it implies, seven populations,
same two architectures throughout.}
\label{tab:dispersion}
\begin{tabular}{lrrlrrr}
\toprule
Population & $n$ & ratio & sd(log ratio) & $95\%$ CI on sd
  & $n$ for $10\%$ & for $20\%$ \\
\midrule
Sine                   & $10$ & $1.184$ & $\mathbf{0.477}$ & $[0.264, 0.590]$ & $\mathbf{150}$ & $40$ \\
Mackey--Glass          & $25$ & $0.907$ & $\mathbf{0.404}$ & $[0.285, 0.501]$ & $\mathbf{150}$ & $30$ \\
Ising 2-spin, $z$      & $10$ & $0.969$ & $0.088$ & $[0.037, 0.120]$ & $8$ & $6$ \\
Ising 4-spin, $zx$     & $10$ & $1.015$ & $0.051$ & $[0.026, 0.062]$ & $6$ & $6$ \\
Ising 4-spin, $z$      & $10$ & $1.001$ & $0.039$ & $[0.022, 0.048]$ & $6$ & $6$ \\
Ising 6-spin, $z$      & $10$ & $0.976$ & $0.030$ & $[0.019, 0.037]$ & $6$ & $6$ \\
Noisy sine             & $10$ & $0.998$ & $\mathbf{0.014}$ & $[0.007, 0.018]$ & $\mathbf{6}$ & $6$ \\
\bottomrule
\end{tabular}
\end{table}

Dispersion varies \textbf{twenty-nine-fold} across populations evaluated
in the same study with the same two architectures ($0.014$ to $0.404$,
or thirty-four-fold including sine), and the intervals for the
forecasting and Ising groups do not overlap. The sample required to
detect a $10\%$ effect varies twenty-fivefold, from six seeds to a
hundred and fifty.

The two extremes deserve note because they are not the tasks one would
guess. The tightest population is \emph{noisy} sine, where added
observation noise evidently stabilises the comparison rather than
destabilising it; the widest is clean sine, which looks like the easiest
task in the study. Dispersion is not predictable from a task's apparent
difficulty, which is the practical reason it has to be measured rather
than assumed.

\subsection{What a study would report when nothing is there}
\label{sec:nullrates}
\figorplaceholder{fig4_power.png}{0.95}{Resampled from the
Mackey--Glass population. A: the estimate a study would report, $90\%$
interval, against seeds run. B: with no true effect, the fraction of
studies reporting an advantage of a given size. C: detection rate at
$p < 0.05$ for true effects of $10\%$, $20\%$ and $30\%$.}{power}

The null is constructed by randomising the sign of each log-ratio ---
exactly the exchangeability the paired test assumes. It preserves the
real dispersion of these models and removes only the direction, so no
distributional assumption is required.

\begin{table}[htb]\centering\small
\caption{What a study with no true effect reports. Resampled under the
paired null from three populations of this study, spanning the
dispersion range of Table~\ref{tab:dispersion}.}
\label{tab:nullrates}
\begin{tabular}{lrrr rrr}
\toprule
& \multicolumn{3}{c}{reports $\ge 10\%$ gain}
& \multicolumn{3}{c}{reaches $p < 0.05$} \\
\cmidrule(lr){2-4}\cmidrule(l){5-7}
Seeds & Sine & Mackey--Glass & Noisy sine
      & Sine & Mackey--Glass & Noisy sine \\
 & \footnotesize sd $0.477$ & \footnotesize sd $0.404$ & \footnotesize sd $0.014$
 & & & \\
\midrule
$5$  & $30.0\%$ & $26.6\%$ & $0.0\%$ & $0.0\%$ & $0.0\%$ & $0.0\%$ \\
$10$ & $23.7\%$ & $19.1\%$ & $0.0\%$ & $4.6\%$ & $3.9\%$ & $5.1\%$ \\
$13$ & $22.0\%$ & $17.0\%$ & $0.0\%$ & $4.6\%$ & $4.7\%$ & $4.3\%$ \\
$20$ & $14.8\%$ & $11.3\%$ & $0.0\%$ & $4.3\%$ & $4.2\%$ & $4.1\%$ \\
\bottomrule
\end{tabular}
\end{table}

The spread across populations is the point. At five seeds on clean sine,
\textbf{three null studies in ten} report a $10\%$-or-better advantage;
on noisy sine, with the same two architectures and the same code,
\textbf{none do}. At thirteen seeds --- above the field's typical sample
--- the sine rate is still better than one in five while the noisy-sine
rate remains zero.

A reader given only a point estimate and a sample size cannot tell these
situations apart. That is the whole difficulty, and it is why the
dispersion of Table~\ref{tab:dispersion} is not an optional
supplementary statistic.

The significance column sits at the nominal $5\%$, confirming the
permutation test is correctly calibrated. The problem is therefore not
false significance. It is that an advantage can be reported as a point
estimate with no test at all, which is common, and under that convention
sampling noise manufactures advantages at a high rate.

This also explains Section~\ref{sec:replication}. Mackey--Glass requires
$30$ seeds to resolve a $20\%$ effect; we ran $13$. The apparent
advantage appeared precisely where measurement was impossible. The Ising
tasks need six seeds for a $10\%$ effect, so their ten-seed nulls are
genuinely informative --- and the three-seed observables contrast failed
not because the task is noisy but because a $5.6\%$ effect needs more
than six seeds even at that low dispersion.

\subsection{A pilot cannot tell you how many seeds you need}
\label{sec:pilot}
The obvious recommendation --- run a small pilot, measure dispersion,
compute the required sample --- does not survive contact with the data.
Treating our $25$-seed population as ground truth and resampling pilots
from it:

\begin{table}[htb]\centering\small
\caption{What a pilot implies about the sample it is meant to forecast.}
\begin{tabular}{lll}
\toprule
Pilot size & sd estimate ($90\%$ range) & implied $n$ for a $10\%$ effect \\
\midrule
$3$  & $0.108$ -- $0.699$ & $\mathbf{4}$ -- $\mathbf{169}$ \\
$5$  & $0.188$ -- $0.594$ & $12$ -- $122$ \\
$8$  & $0.244$ -- $0.548$ & $21$ -- $104$ \\
$10$ & $0.270$ -- $0.523$ & $25$ -- $95$ \\
$13$ & $0.300$ -- $0.490$ & $31$ -- $83$ \\
\bottomrule
\end{tabular}
\end{table}

A three-seed pilot implies a required sample anywhere between four and a
hundred and sixty-nine. Even a ten-seed pilot leaves a four-fold range.
Dispersion is a second-moment quantity and is estimated poorly at small
$n$ --- so the practical recommendation cannot be ``forecast your sample
size'', and must instead be ``report the interval, which encodes the
resolving power you actually achieved''.

\subsection{How much of ``seed variance'' is seed variance?}
\label{sec:repeat}
The protocol of Section~\ref{sec:setting} holds the data split fixed so
that seeds vary model initialisation only. That is the intent; whether
the pipeline honours it is a separate question, and we can answer it
because three seeds of one forecasting task were inadvertently run
twice at identical seed and configuration.

\begin{table}[htb]\centering\small
\caption{Two runs at identical seed and configuration, clean sine.
Absolute test error and the paired quantum/classical ratio.}
\label{tab:repeat}
\begin{tabular}{lrrr}
\toprule
& seed $0$ & seed $1$ & seed $2$ \\
\midrule
quantum, run A   & $2.63\times10^{-5}$ & $2.34\times10^{-5}$ & $2.77\times10^{-5}$ \\
quantum, run B   & $3.30\times10^{-5}$ & $4.59\times10^{-5}$ & $2.09\times10^{-5}$ \\
classical, run A & $1.29\times10^{-5}$ & $2.14\times10^{-5}$ & $3.36\times10^{-5}$ \\
classical, run B & $1.56\times10^{-5}$ & $4.37\times10^{-5}$ & $2.28\times10^{-5}$ \\
\midrule
ratio, run A     & $2.032$ & $1.092$ & $0.823$ \\
ratio, run B     & $2.112$ & $1.050$ & $0.913$ \\
\bottomrule
\end{tabular}
\end{table}

\textbf{Absolute error moves by up to a factor of two at a fixed seed.}
Seed control does not reproduce in this pipeline, and the same is true
in the companion study, where one closed-loop run differed by a factor
of $3.6$ between two executions at nominally identical seed and
configuration~\cite{paperb}. We have not established the cause; the
leading hypothesis is that thread count alters floating-point reduction
order and so the optimisation trajectory.

\textbf{The paired ratio is far more stable.} The log standard deviation
of the run-to-run change in the ratio is $0.071$, against $0.477$ across
seeds on the same task --- a factor of seven. Both arms share the data
and most of the pipeline, so whatever perturbs one perturbs the other,
and the pairing absorbs it.

Two consequences. The paired design is justified empirically here rather
than by assertion, which we had not previously been able to claim. And
a study reporting \emph{unpaired} comparisons inherits the full
run-to-run component on top of the seed component --- which is one
mechanism by which the same thirteen seeds in
Section~\ref{sec:replication} give $p = 0.022$ unpaired and $p = 0.069$
paired.

Six slots is a thin basis and we report it as indicative. It is,
however, a direct measurement of a quantity usually assumed to be zero.

% =====================================================================
\section{The permutation floor}\label{sec:floor}
% =====================================================================
One constraint in this area is arithmetic rather than statistical, and
it is violated often enough to be worth stating separately.

An exact two-sided sign-flip permutation test on $n$ paired
observations enumerates $2^n$ sign assignments. The observed assignment
and its exact negation are always at least as extreme as themselves, so
the smallest two-sided $p$-value the test can return is
\begin{equation}
  p_{\min}(n) = \frac{2}{2^{n}} = 2^{1-n}.
\end{equation}

\begin{table}[htb]\centering\small
\caption{The smallest attainable two-sided $p$-value, by sample size.}
\label{tab:floor}
\begin{tabular}{lrrrrrrr}
\toprule
$n$ & $3$ & $4$ & $5$ & $6$ & $7$ & $8$ & $10$ \\
\midrule
$p_{\min}$ & $0.250$ & $0.125$ & $0.062$ & $0.031$ & $0.016$ & $0.008$ & $0.002$ \\
\bottomrule
\end{tabular}
\end{table}

A three-seed comparison therefore \emph{cannot} reach $p < 0.05$, and a
five-seed comparison cannot either, whatever the data and however large
the effect. This is independent of dispersion: it is a property of the
test's support. Four of the eight comparisons in
Table~\ref{tab:forest} sit at three seeds, and we report their
$p$-values as $0.500$, $0.750$, $0.500$ and $1.000$ precisely because
nothing smaller was available.

The practical consequence is that a study reporting ``ahead in all three
seeds'' is reporting the most extreme outcome its design admits, which
is not the same as evidence. The floor should be stated alongside any
small-sample $p$-value, and it costs nothing to compute.

\paragraph{Multiplicity, across both studies.}
Reporting many comparisons raises the complementary problem. Taking
every paired test in this paper and in the companion architecture
study~\cite{paperb} together:

\begin{table}[htb]\centering\small
\caption{Nominally significant results against what the global null
predicts.}
\label{tab:multiplicity}
\begin{tabular}{lrrr}
\toprule
& This paper & Companion study & Both \\
\midrule
Paired comparisons reported & $7$ & $22$ & $30$ \\
Reaching $p < 0.05$         & $1$ & $2$  & $3$ \\
Expected under the null     & $0.35$ & $1.1$ & $1.5$ \\
Bonferroni threshold        & $0.0071$ & $0.0023$ & $0.0017$ \\
Smallest $p$ observed       & $0.041$ & $0.039$ & $0.039$ \\
\bottomrule
\end{tabular}
\end{table}

Three nominally significant cells among thirty comparisons, against
$1.5$ expected, is close to what one observes when there is no effect
anywhere, and none of the three approaches the corrected threshold. The
three are the Ising six-spin comparison of Table~\ref{tab:forest}
($p = 0.041$), an offline cell in the companion study ($0.041$), and
that study's closed-loop result ($0.039$).

We draw the conclusion for our own work before drawing it for anyone
else's: \textbf{none of our three $p < 0.05$ results should be treated
as established without replication at fresh seeds.} The companion study
says so of its headline, and we say so here of the six-spin cell. This
is uncomfortable --- they are the only positive findings either paper
has --- and it is what our own argument requires.

% =====================================================================
\section{Silent failures: a case series}\label{sec:silent}
% =====================================================================
The preceding sections concern whether a measurement can resolve its
target. This section concerns whether the pipeline that produced the
measurement did what its author believed. Across this study and the
companion architecture study~\cite{paperb}, seven failures occurred that
produced plausible wrong numbers rather than errors. Four changed a
reported result. We report them because the methodological literature in
this area is almost entirely statistical, and in our experience these
were the failures that mattered.

Five were found in the companion study. \textbf{Two were found in this
paper, by running on our own pipeline the audit this section
recommends}, after it had been written. We report those in
Section~\ref{sec:selfaudit} rather than burying them, because a paper
arguing that aggregation errors change published numbers has no standing
to exempt itself.

All five share one structure: \textbf{a quantity that varied was not
recorded as having varied.}

\subsection{The instructive one: a $p$-value that did not move}
\label{sec:pooling}
An aggregation key grouped runs by model, capacity and horizon, but not
by \emph{task}. Runs from two different environments --- a classical
control task and a pendulum --- therefore collided on shared seeds, and
because one task name sorts after the other, one value silently replaced
its counterpart at every shared key.

The two environments have very different error scales: classical
one-step error is approximately $0.00107$ on the first and $0.00425$ on
the second. A single pooled run therefore inflated the classical arm's
mean fourfold, manufacturing an apparent quantum improvement.

\begin{table}[htb]\centering\small
\caption{The same comparison before and after the task-pooling error
was found. Eight seeds throughout; the sign pattern never changed.}
\label{tab:pooling}
\begin{tabular}{lrrr}
\toprule
& ratio & $p$ & quantum ahead \\
\midrule
As first computed (two tasks pooled) & $\mathbf{0.711}$ & $0.041$ & $6/8$ \\
After the key was corrected          & $\mathbf{0.975}$ & $0.041$ & $6/8$ \\
\bottomrule
\end{tabular}
\end{table}

A $29\%$ improvement becomes $2.5\%$. \textbf{The $p$-value is identical
to three decimal places}, because the permutation test operates on the
paired differences and the sign pattern was unaffected; only the
magnitude was spurious. The result was significant at $\alpha = 0.05$
before the error and significant after it, and a reader applying any
standard statistical check would have found nothing wrong.

We emphasise this because it bounds what statistical practice can
achieve. Reporting intervals, as this paper recommends throughout, would
not have caught it either: the interval would have been computed from
the same contaminated population and would have excluded unity just as
confidently. The error was found only because the aggregation code was
made to halt when two runs collide on a key.

\subsection{The other four}
\begin{enumerate}\itemsep4pt
\item \textbf{A run-naming collision destroyed two completed runs.} A
      directory-naming scheme appended only non-default settings, so an
      ablation at a default value produced the same directory name as a
      headline run at a different step budget and overwrote it. Both
      runs completed normally and left plausible result files. The loss
      was invisible until file timestamps were compared against the
      numbers. Cost: two runs of roughly twenty-seven hours each.
\item \textbf{A transformation mismatch inverted a comparison.} An
      open-loop prediction measurement compared decoder output against
      raw observations, but the decoder is trained against
      $\mathrm{symlog}$ of the observation. The mismatch applied an
      unmodelled transformation to one side only, inflating the model's
      error while leaving the persistence baseline untouched, and
      reversed the sign of the reported result.
\item \textbf{A confounded ablation.} A sweep over attention context
      length varied model width simultaneously --- context $2$ at width
      $32$, context $4$ at width $8$, context $8$ at width $6$ --- so
      the three points were not comparable and no trend could be read
      from them. The naming guard introduced after item~1 prevents
      overwriting but not confounding, so a separate check was added
      that refuses to start a run when a result-affecting setting absent
      from the directory name differs from its default.
\item \textbf{An analysis grouping merged distinct conditions.} A paired
      comparison keyed on seed and arm alone, so every experimental
      condition collided and whichever directory sorted last silently
      won. It happened to select the intended condition, by alphabetical
      accident rather than by construction.
\end{enumerate}

\subsection{Two in this paper, found by the recommended audit}
\label{sec:selfaudit}
We wrote an audit script that rebuilds each published number from the
raw records under a full-configuration key and halts on collision, and
ran it against this paper. It failed two checks.

\paragraph{A filter that never filtered.}
The loader supplying the latent-dynamics row of Table~\ref{tab:forest}
selected rows for which a width field was absent, intending to pick the
main experimental arm. That field is absent from \emph{every} record in
the file, so the filter admitted all of them --- including five
different model widths from a capacity ablation --- and the subsequent
key, on seed and model only, retained whichever was read last. Since
Section~\ref{sec:sensitivity} establishes that width alone moves error
by $13\%$, the row was not a measurement. It is withdrawn above.

\paragraph{A data file containing three experiments and a wrong
entangler.}
The forecasting records are a single file that had been appended to
three times. It contains a failed run (test error of order $0.8$ on a
sine wave, worse than predicting zero), two repeats of seeds $0$--$2$ at
identical configuration, and seeds $3$--$9$ carrying an eighth field
--- \texttt{ring\_cascade} --- that the file's own header does not
declare. The loader, keyed on signal, seed and arm, retained the last
matching row in each case.

Two consequences. The failed run was excluded by ordering rather than by
intent, which is luck. And the entangler is a ring cascade, not the
brickwork this paper's methods section claimed until
Section~\ref{sec:setting} was corrected --- a discrepancy of some
embarrassment in a paper whose Section~\ref{sec:trainability} measures a
$1{,}260\times$ consequence of exactly that choice.

The one benefit: the duplicate repeats are what made
Section~\ref{sec:repeat} possible. We would not otherwise have been able
to measure the fixed-seed reproducibility of our own pipeline.

\paragraph{Why these were not caught sooner.}
Both loaders produce numbers of entirely plausible magnitude, and both
had been read by their author several times. Neither raises an error.
The published ratios they produced --- $0.991$ for the pooled widths,
and forecasting figures that are individually unremarkable --- carry no
signature of the defect. We note that the audit which found them is
about a hundred lines and took under a minute to run.

\subsection{What we conclude}
Four of the seven changed a reported number. Four were found only
because a check was added that had not previously existed, and we have
no reason to think the list is complete --- which is itself the point.
The two most recent were found in the paper you are reading, by the
procedure that paper recommends, after the recommendation had been
drafted.

Two lessons, stated as consequences rather than advice.

\textbf{An aggregation key written by enumerating the fields one is
thinking about will omit the fields one is not.} The correct default is
to key on the entire configuration and exclude only the axes under
comparison, and to halt rather than select a winner when two runs
collide. Every aggregation in this paper and in~\cite{paperb} now does
so. This check costs a few lines and found two errors that had already
propagated into draft text.

\textbf{Run provenance deserves the same standing as interval
estimation.} The field's methodological discussion treats statistics as
the locus of error. In this project the statistics were calibrated
throughout --- the permutation test sits at its nominal rate in
Section~\ref{sec:nullrates} --- and the errors that changed our
conclusions were all failures of bookkeeping. Section~\ref{sec:pooling}
is the sharp case: a significance test passing a fabricated effect size
without flinching.

\pendpara{We would like to report whether this class of error is common
in published work, but it is not detectable from a paper --- only from
its pipeline. We note that as a limitation of the present contribution
and as an argument for releasing aggregation code rather than only
results.}

% =====================================================================
\section{A trainability diagnostic that measures the wrong variable}
\label{sec:trainability}
% =====================================================================
\figorplaceholder{fig2_entangler.png}{0.95}{Gradient variance against
circuit width, by cost-function locality and entangler topology, at two
entangling layers.}{entangler}

Cerezo et al.~\cite{cerezo2021} establish that local cost functions
avoid barren plateaus when ansatz depth is $O(\log n)$. The condition is
on \emph{depth}, via the causal cones of the cost terms. We report what
happens when a common implementation choice violates it silently.

\begin{table}[htb]\centering\small
\caption{Decay factor in gradient variance per added qubit, log-linear
fit over $2$--$10$ qubits.}
\begin{tabular}{llrrr}
\toprule
Cost & Entangler & 2 layers & 4 layers & 6 layers \\
\midrule
global & brickwork    & $\times 0.366$ & $\times 0.484$ & $\times 0.500$ \\
global & ring cascade & $\times 0.464$ & $\times 0.484$ & $\times 0.496$ \\
local  & brickwork    & $\times 0.988$ & $\times 0.963$ & $\times 0.842$ \\
local  & ring cascade & $\times 0.491$ & $\times 0.501$ & $\times 0.504$ \\
\bottomrule
\end{tabular}
\end{table}

Under a local cost at two entangling layers, brickwork shows essentially
no decay ($0.2206 \rightarrow 0.2161$ across two to ten qubits) while a
sequential ring cascade decays by a factor of $295$ over the same range
--- roughly $\mathbf{1{,}260\times}$ apart at ten qubits. The
global-cost rows are the control: there the topologies are
indistinguishable, as the theory predicts, since global costs are
unprotected at any depth.

The explanation is not novel; it is~\cite{cerezo2021} applied
correctly. A ring cascade applies CNOTs sequentially, so correlations
traverse the register within one nominal layer and the causal cone is
immediately global. A brickwork layer is depth $2$ at any width.
Plotting gradient variance against \emph{layers} with a sequential
entangler therefore measures depth, not layer count, and the
shallow-depth precondition fails while the axis label still reads two
layers.

The result is a practitioner one with a measured magnitude. An
implementation using a cascade entangler and plotting against layers
will exhibit a plateau the theory does not predict for that cost
function, inviting attribution to the phenomenon rather than to the
circuit. The magnitude matters because this diagnostic is used
prospectively, to choose qubit counts and circuit depths before the main
experiments are run --- as it was here. A confounded diagnostic
therefore propagates into design decisions, not only into conclusions,
and does so at a point where the error is no longer visible.

\paragraph{The measurement cost this implies.}
Because shot noise on an estimated expectation scales as $1/\sqrt{N}$,
the shots needed to resolve a gradient of magnitude $g$ scale as
$g^{-2}$. At ten qubits with local readout, resolving a gradient to
fixed relative precision requires approximately $3$ shots with
brickwork and $9{,}360$ with a ring cascade --- a factor of
$\mathbf{2{,}766}$. For simulation this is a stylistic difference; for
hardware it is the difference between a feasible experiment and an
impossible one.

These estimates use $100$ gradient samples per configuration, which is
thin. \pend{rerun at $1{,}000$ samples; roughly one hour}

% =====================================================================
\section{Why no effect is expected at this scale}\label{sec:why}
% =====================================================================
The results above bound an effect without explaining its absence. A
bound without a mechanism invites the obvious objection --- that a
better circuit would have worked --- and that objection deserves a
direct answer.

\subsection{The circuit computes a truncated Fourier series}
Schuld, Sweke and Meyer~\cite{schuld2021} show that a variational model
with data re-uploading computes a partial Fourier series in its inputs,
in which the accessible frequencies are fixed by the encoding gates and
the trainable parameters set only the coefficients. The circuit used
here encodes each input dimension through an $R_X$ and an $R_Z$ rotation
in every layer, so with $L$ layers each dimension passes through $2L$
encoding gates and its accessible frequencies are integers bounded by
that count.

The consequence is that the quantum block is not a richer function class
than the classical one it replaces. At this scale it is a \emph{more
constrained} one: a trigonometric feature map over a fixed integer
spectrum, with $28$ parameters setting the coefficients subject to
whatever constraints the circuit structure imposes, at roughly
$20\times$ the computational cost. A classical model of comparable
budget over the same spectrum spans a superset of it.

This also accounts for Section~\ref{sec:sensitivity}. Classical width
buys $13\%$ because capacity binds on that task; the quantum
substitution buys $-0.1\%$ because it does not add capacity, but
exchanges one small parameterisation for another.

\subsection{The only remaining route is inductive bias}
Expressivity advantage is unavailable by construction wherever the
circuit is classically simulable: any function it computes can be
computed classically at comparable cost. What remains is the possibility
that the circuit's constrained parameterisation constitutes a useful
prior, in the way convolutional structure does for images. Our Ising
measurement bounds any such effect at $2.1\%$.

\subsection{A classical surrogate, and a prior question it exposes}
\label{sec:surrogate}
We construct the classical counterpart directly: identical down- and
up-projections, with the circuit replaced by a map over fixed integer
frequencies drawn from its accessible spectrum and learned
coefficients. The coefficient budget is swept to locate the classical
parameter count at which the surrogate matches the circuit.

The spectrum at four input dimensions and six layers contains on the
order of $10^5$ frequency vectors, too many to span, so the surrogate
samples it in the manner of random Fourier features. The evidence is
therefore asymmetric: a sample that suffices is strong evidence that
nothing distinctive is being exploited, whereas a sample that does not
suffice is equally consistent with a poor draw.

\figorplaceholder{fig5_surrogate.png}{0.95}{Chaotic forecasting, five
seeds. Left: test error against trainable parameters in the
token-mixing block; the shaded band is the circuit's own $10$--$90\%$
range across seeds. Right: each arm paired against the circuit, with
$95\%$ bootstrap intervals.}{surrogate}

\begin{table}[htb]\centering\small
\caption{Classical surrogate sweep against the quantum circuit.}
\label{tab:surrogate}
\begin{tabular}{lrrlr}
\toprule
Arm & block parameters & ratio to circuit & $95\%$ CI & $p$ \\
\midrule
Fourier surrogate & $8$   & $\mathbf{1.029}$ & $[0.890, 1.210]$ & $0.750$ \\
Fourier surrogate & $16$  & $0.861$ & $[0.657, 1.092]$ & $0.438$ \\
Fourier surrogate & $32$  & $1.044$ & $[0.791, 1.325]$ & $0.750$ \\
Fourier surrogate & $64$  & $1.093$ & $[0.873, 1.318]$ & $0.500$ \\
Fourier surrogate & $128$ & $0.918$ & $[0.689, 1.210]$ & $0.625$ \\
Fourier surrogate & $256$ & $1.121$ & $[0.893, 1.342]$ & $0.375$ \\
Fourier surrogate & $512$ & $1.164$ & $[0.900, 1.506]$ & $0.375$ \\
Classical bottleneck & $460$ & $1.284$ & $[0.948, 1.739]$ & $0.250$ \\
\bottomrule
\end{tabular}
\end{table}

The strong direction fires. \textbf{A classical Fourier map with eight
coefficients --- a single sampled frequency vector --- matches the
thirty-parameter circuit} (ratio $1.029$, interval $[0.890, 1.210]$). A
random sample of one frequency from a spectrum of $10^5$ suffices, which
is difficult to reconcile with the circuit exploiting anything
distinctive about its function class.

\paragraph{But the sweep reveals why, and the reason is more
consequential.}
Varying the block's capacity over a $64$-fold range, from $8$ to $512$
parameters, barely moves the error: the log--log slope is $+0.038$ and
the best configuration is only $0.741$ of the worst.

\begin{table}[htb]\centering\small
\caption{Where the variation actually is.}
\begin{tabular}{lr}
\toprule
Quantity & sd of log error \\
\midrule
Across the entire $64\times$ capacity sweep & $\mathbf{0.109}$ \\
Across random seeds, quantum circuit alone  & $\mathbf{0.243}$ \\
\bottomrule
\end{tabular}
\end{table}

\textbf{Changing the block's capacity by a factor of $64$ moves
performance less than half as much as changing the random seed.} The
token-mixing block does not constrain performance on this task. Whatever
is placed in that slot --- a quantum circuit, an eight-parameter
trigonometric map, a $460$-parameter classical bottleneck --- performs
identically, because the block is not where the task's difficulty lies.

This subsumes the Fourier question on this benchmark rather than
answering it. The surrogate matches the circuit, but so would almost
anything, so the match cannot be read as evidence specifically about
function classes. The claim we can make is narrower and firmer:
\textbf{on this task, no comparison sited at the token-mixing block is
informative, at any sample size.}
\pend{run the surrogate sweep on the Ising task, where
Section~\ref{sec:sensitivity} establishes that capacity binds at the
model level, though not specifically at this block}

\subsection{An independent test: remove the entanglement}
\label{sec:entanglement}
The companion study~\cite{paperb} provides a second, independent probe
of the same question on a different architecture. Every two-qubit gate
is deleted from the circuit, so it becomes a product of independent
single-qubit rotations and can generate no entanglement whatever; the
parameter count is unchanged, so the comparison is otherwise exact.

\begin{table}[htb]\centering\small
\caption{Entanglement ablation, from~\cite{paperb}. Eight seeds, six
qubits. Both arms are quantum and differ only in the presence of
two-qubit gates.}
\begin{tabular}{lrr}
\toprule
Entangler & 1-step error & 10-step error \\
\midrule
Brickwork & $0.00106$ & $0.00530$ \\
None      & $0.00107$ & $0.00536$ \\
\midrule
Ratio (none/brickwork) & $1.009$ & $1.011$ \\
\bottomrule
\end{tabular}
\end{table}

Removing all entanglement changes error by about $1\%$. This is stronger
than the Fourier argument alone. That argument establishes the
function class is classically reachable \emph{in principle}; the
ablation shows the property that would make the circuit hard to simulate
is contributing nothing \emph{in practice}. Two independent routes --- one
theoretical, one measured --- reach the same place, and the second
corroborates the entanglement finding of~\cite{bowles2024} on an
architecture they did not test.

\subsection{A second necessary condition}\label{sec:second}
Section~\ref{sec:surrogate} generalises beyond the present study, and it
is independent of everything in Section~\ref{sec:howmany}. A comparison
between architectural components requires two conditions, and failing
either makes the result uninformative:

\begin{enumerate}\itemsep2pt
\item \textbf{The sample must resolve the effect} --- the dispersion
      requirement of Section~\ref{sec:howmany}, subject to the floor of
      Section~\ref{sec:floor}.
\item \textbf{The component under comparison must constrain
      performance} --- otherwise every candidate performs alike, and the
      comparison measures nothing whatever the sample size.
\end{enumerate}

The second is cheaper to check than the first: sweeping the component's
capacity costs a handful of classical runs, needs no quantum simulation,
and returns a direct answer. Its failure mode is also more insidious,
because adding seeds cannot fix it --- a well-powered comparison at a
position that does not bind yields a precise measurement of zero, easily
misread as a meaningful null.

The architecture evaluated here, and the benchmark on which its
advantage was originally reported, fail this condition. That offers a
simpler account of why a capacity-matched classical bottleneck matched
the circuit in~\cite{qasa} than the architectural-parsimony explanation
those authors give: at this position on this task, matching requires no
particular architecture at all.

\subsection{The scale window}
Taken together with the trainability literature, this suggests a
structural difficulty extending beyond the present study. Below the
simulability threshold, expressivity advantage is impossible by
construction and only inductive-bias effects remain available. Above it,
three obstacles arrive at once: gradient magnitudes suppressed with
system size~\cite{cerezo2021}; the unavailability of backpropagation,
since parameter-shift requires two shot-limited circuit evaluations per
parameter; and, for quantum-native data specifically, the fact that
generating a training set requires the simulation presumed infeasible.

The regime in which an advantage is both possible in principle and
trainable in practice is therefore narrow, and may with current methods
be empty. We do not claim it is: we observe that the present
experimental paradigm cannot reach it, which is a constraint on what the
field's empirical results can currently establish.

\subsection{What this does not rule out}
Every circuit evaluated here belongs to the hardware-efficient family
--- problem-agnostic ans\"atze of fixed structure. Two directions remain
untested. A \textbf{problem-informed ansatz}, constructed from the terms
of the Hamiltonian being modelled, imposes structure the generic family
does not and would alter the accessible spectrum in a task-dependent
way. \textbf{Non-variational approaches} avoid the optimisation
obstacles entirely: Pandey~\cite{pandey2026} reports a fixed quantum
reservoir outperforming a variational quantum model by $81\%$ on chaotic
prediction while training some four orders of magnitude faster ---
evidence that the limiting factor may be the variational paradigm rather
than the quantum substrate.

% =====================================================================
\section{What follows}\label{sec:follows}
% =====================================================================
Five consequences follow from the measurements above. We state them as
implications of the data rather than as procedural advice, since each is
quantified by a specific result.

\paragraph{Dispersion determines interpretability, and must be
reported.}
Section~\ref{sec:dispersion} shows the same architecture pair exhibiting
dispersion from $0.039$ to $0.477$ across settings, and the sample
required to detect a $10\%$ effect moving from $6$ to $150$ accordingly.
A reported sample size is uninterpretable without it: ten seeds is ample
in one of our settings and an eighth of what is needed in another. The
quantity is a by-product of runs already performed.

\paragraph{Sample size cannot be forecast from a pilot, so intervals
must carry the argument.}
Section~\ref{sec:pilot} shows a three-seed pilot implying a required
sample between $4$ and $170$. Prospective sample-size calculation is not
available in this regime, and the confidence interval --- which encodes
resolving power achieved rather than predicted --- has to do the work
instead. A $2.1\%$ bound and a $22\%$ bound appear identical in a table
of point estimates.

\paragraph{Point estimates without intervals convert dispersion into
apparent advantages.}
Section~\ref{sec:nullrates} quantifies the rate: under the paired null
at the dispersion our highest-variance task exhibits, $29\%$ of
five-seed studies and $17\%$ of thirteen-seed studies report a
$10\%$-or-better advantage. The permutation test remains correctly
calibrated throughout, so this is not a failure of significance testing
but a consequence of reporting effects without one.

\paragraph{Small-sample $p$-values should be reported with their floor.}
Section~\ref{sec:floor}: at three seeds no exact paired test can return
$p < 0.25$, and at five seeds none can return $p < 0.062$. A study
reporting ``ahead in all three seeds'' is reporting the most extreme
outcome its design admits. The floor is one line of arithmetic.

\paragraph{Aggregation code deserves the scrutiny given to statistics.}
Section~\ref{sec:silent}: three of five silent failures in the companion
study changed a reported number, and the most instructive left the
$p$-value unchanged while fabricating the effect size. Keying
aggregations on the complete configuration, and halting on collision
rather than selecting a winner, is a few lines of code and caught two
errors that had already reached draft text. We recommend releasing
aggregation code alongside results, because this class of error is
invisible from a paper.

% =====================================================================
\section{Limitations}
% =====================================================================
\begin{itemize}\itemsep3pt
\item A single architecture family. Every circuit evaluated is
      hardware-efficient and variational. The Fourier argument is
      specific to data re-uploading with Pauli encodings and does not
      constrain problem-informed ans\"atze or non-variational schemes.
\item Exact expectation values throughout. Finite-shot behaviour, and so
      all behaviour on hardware, is largely untested here. As argued in
      Section~\ref{sec:setting} this makes the null conservative rather
      than fragile, but it leaves the hardware regime unexamined.
      \pend{the shot-noise probe exists; either run it or remove the
      claim}
\item Baselines matched by parameter count rather than compute, which
      favours the quantum model given its $\sim\!20\times$ wall-clock
      cost.
\item Wall-clock measurements contaminated by machine sleep and
      competing load.
\item Gradient-variance measurements at $100$ samples per configuration.
\item Several comparisons remain at three seeds and are reported as
      uninformative rather than as evidence of parity.
\item The dispersion measurement rests on seven populations. The
      twenty-nine-fold range is a lower bound on the variation that
      exists across tasks generally; we cannot say where the extremes
      lie.
\item The forecasting experiments use a ring-cascade entangler and the
      remainder use brickwork (Section~\ref{sec:setting}). None of our
      conclusions depends on the choice, but the study is not uniform
      in it.
\item The fixed-seed reproducibility measurement of
      Section~\ref{sec:repeat} rests on six slots and is indicative
      only.
\item One comparison (Ising, six spins) reaches $p < 0.05$ and has not
      been replicated at fresh seeds. By this paper's own argument it
      should not be treated as established.
\item Dispersion is measured for one architecture pair. Whether it is a
      property of the task, the architecture, or their interaction is
      not determined here, and that question bears directly on how
      transferable the recommendation is.
\item The case series of Section~\ref{sec:silent} is from one project by
      one author. It establishes that this class of failure occurs and
      what it costs; it establishes nothing about its prevalence.
\end{itemize}

% =====================================================================
\section*{Code and data availability}
% =====================================================================
All code, run records and analysis scripts at \pend{repository URL and
archival DOI}. Every figure and table is produced from the raw run
records by \texttt{make\_figures.py} and \texttt{power\_analysis.py};
intermediate values are written to \texttt{figures/figure\_data.json}
and \texttt{figures/power\_data.json} so that numbers in the text can be
checked against numbers in the plots. Aggregation keys and their
collision checks are in the same scripts, for the reason given in
Section~\ref{sec:silent}.

% =====================================================================
\section*{Acknowledgements}
% =====================================================================
\pend{acknowledgements, funding, hardware access}

% =====================================================================
\appendix
\section{Remaining work}\label{app:remaining}
% =====================================================================
\begin{enumerate}\itemsep3pt
\item \textbf{Done.} Regenerated Table~\ref{tab:forest} and
      Table~\ref{tab:dispersion}; the dispersion measurement now rests
      on seven populations rather than four.
\item \textbf{Done.} Audit of every number against the raw run records
      (\texttt{audit\_paperA.py}). It found the two failures reported in
      Section~\ref{sec:selfaudit}.
\item \textbf{Blocking.} Repair the forecasting data file: record the
      entangler explicitly, remove the failed run, and retain the
      duplicate repeats as a labelled subset rather than letting row
      order select among them. Then rebuild
      \texttt{make\_figures.py}'s loaders to key on the full
      configuration and halt on collision, as the other analysis scripts
      now do.
\item \textbf{Blocking.} Reinstate the latent-dynamics row from the
      Phase 1 records at eight seeds, with the context length stated.
\item Replicate the Ising six-spin comparison at fresh seeds. It is the
      only cell in Table~\ref{tab:forest} below $\alpha = 0.05$ and
      Section~\ref{sec:floor} requires us to treat it as provisional.
\item Exact citations for
      \cite{henderson2018,agarwal2021,bouthillier2021,dodge2019,pandey2026};
      author list for \cite{qasa}; journal details for
      \cite{schuld2021}.
\item Rerun gradient variance at $1{,}000$ samples
      (Section~\ref{sec:trainability}).
\item Bring the noisy-sine setting to ten seeds.
\item Run the surrogate sweep on the Ising task, where capacity binds
      (Section~\ref{sec:surrogate}).
\item Run the shot-noise probe, or remove the finite-shot claim.
\item Rerun the classical width sweep of
      Section~\ref{sec:sensitivity} at three seeds so the table can
      carry intervals.
\item Search for prior art: any QML paper reporting seed dispersion and
      deriving a required sample size; any first-hand account of
      aggregation-key failure in an ML pipeline.
\item Repository URL and archival DOI.
\end{enumerate}

% =====================================================================
\bibliographystyle{unsrtnat}
\begin{thebibliography}{9}
\bibitem{bowles2024} J.~Bowles, S.~Ahmed, M.~Schuld. \emph{Better than
  classical? The subtle art of benchmarking quantum machine learning
  models.} arXiv:2403.07059 (2024).
\bibitem{qasa} \emph{Quantum Adaptive Self-Attention for Quantum
  Transformer Models.} arXiv:2504.05336 (2025).
  \pend{author list from the revised version}
\bibitem{cerezo2021} M.~Cerezo, A.~Sone, T.~Volkoff, L.~Cincio,
  P.~J.~Coles. \emph{Cost function dependent barren plateaus in shallow
  parametrized quantum circuits.} Nature Communications \textbf{12},
  1791 (2021).
\bibitem{schuld2021} M.~Schuld, R.~Sweke, J.~J.~Meyer. \emph{The effect
  of data encoding on the expressive power of variational quantum
  machine learning models.} Phys.\ Rev.\ A \textbf{103}, 032430 (2021);
  arXiv:2008.08605. \pend{confirm volume and page}
\bibitem{henderson2018} P.~Henderson et al. \emph{Deep Reinforcement
  Learning that Matters.} \pend{full citation}
\bibitem{agarwal2021} R.~Agarwal et al. \emph{Deep Reinforcement
  Learning at the Edge of the Statistical Precipice.} \pend{full
  citation}
\bibitem{bouthillier2021} X.~Bouthillier et al. \emph{Accounting for
  Variance in Machine Learning Benchmarks.} \pend{full citation}
\bibitem{dodge2019} J.~Dodge et al. \pend{full citation; on
  computational cost in model comparison}
\bibitem{pandey2026} T.~Pandey. \emph{Fixed-Reservoir vs Variational
  Quantum Architectures for Chaotic Dynamics: Benchmarking QRC and QPINN
  on the Lorenz System.} arXiv:2604.23743 (2026). \pend{verify}
\bibitem{paperb} V.~Liutarevich. \emph{A Quantum Transformer World Model
  for Model-Based Reinforcement Learning.} \pend{companion paper; cite
  by arXiv identifier once posted}
\end{thebibliography}

\end{document}
